% !PRINTED 309
\weberpage{309}
%<<CONT
so the functions (1), apart from positive numerical factors, take for $x = 0$ the values
\weq{2}{a_n, a_{n-1}, a_{n-2} \ldots a_0,}
and from this, accordingly, we obtain the theorem:

\gesperrt{IV. The number of positive roots of the equation $f(x) = 0$ is equal to, or less by an even number than, the number of changes of sign in the sequence of the coefficients of $f(x)$, where any vanishing coefficients are simply to be passed over, and the root $x = 0$, if it is present, is not counted.}

To estimate the number of negative roots as well, one can either apply the \gesperrt{Budan-Fourier} theorem to the interval $(-\infty, 0)$, or one replaces $x$ by $-x$, that is, one changes the signs of $a_1, a_3, a_5 \ldots$ and then applies Theorem IV.

In a similar way we shall specialise the \gesperrt{Newton} criterion of \S.\,100.

We remark for this purpose that
\[
f_\nu(0) = \frac{\Pi(n-\nu)\,\Pi(\nu)}{\Pi(n)}\, a_{n-\nu},
\]
and consequently
\[
F_\nu(0) = \frac{\Pi(\nu)\,\Pi(\nu-1)\,\Pi(n-\nu)\,\Pi(n-\nu-1)}{\Pi(n)\,\Pi(n)}\left[\nu(n-\nu)\,a_{n-\nu}^2 - (\nu+1)(n-\nu+1)\,a_{n-\nu-1}\,a_{n-\nu+1}\right]
\]
and further that the sequence of the $f_\nu(x)$ presents, for $x = +\infty$, only successions of signs, and for $x = -\infty$ only changes of sign.

If, then, we set
\weq{3}{A_\nu = \nu(n-\nu)\,a_\nu^2 - (\nu+1)(n-\nu+1)\,a_{\nu-1}\,a_{\nu+1},}
\[
A_0 = 1, \quad A_n = 1,
\]
so that $A_\nu$, apart from a positive factor, agrees with $F_{n-\nu}(0)$, it follows that for $x = 0$ the double sequence (7), \S.\,100, written in reverse order, agrees in sign with
\weq{4}{
\begin{array}{c}
a_0,\ a_1,\ a_2 \ldots a_n \\
A_0,\ A_1,\ A_2 \ldots A_n.
\end{array}
}
On passing to $x = +\infty$, in the sequence $f(x), f_1(x) \ldots f_n(x)$, all the changes of sign are, and on passing to
%CONT>>
